\section{Introducción}\label{sec:intro}

El modelo de la entrega predice el rendimiento de cebada con dos datos de cada parcela: el régimen de su etiqueta y
el rendimiento de sus vecinas del mismo régimen. Funciona para las 59 parcelas del reto porque están intercaladas con
las de entrenamiento, pero no responde la pregunta que más le serviría a FIRA: \emph{si se siembra una parcela en un
lugar sin parcelas medidas, ¿qué rendimiento se puede esperar?} Para eso hace falta una variable del lugar (clima,
suelo o terreno) que exista en cualquier punto del mapa.

El régimen es la pista natural. Si los cuatro municipios del régimen A rinden 1.44\,t/ha más que los cinco del B,
algo debería estar más presente en unos que en otros. Este reporte busca ese algo. Cada hipótesis de lugar se somete
a las mismas pruebas, y la Fig.~\ref{fig:embudo} resume el recorrido. Primero medimos si el ambiente generaliza a
municipios nuevos (Sección~\ref{sec:generalizacion}). Después descartamos una por una las causas de lugar
(Sección~\ref{sec:descarte}) y contrastamos el dataset con las cifras oficiales (Sección~\ref{sec:siap}). Al final
separamos lo que podemos afirmar de lo que no (Sección~\ref{sec:discusion}).

\begin{figure*}[t]
\centering
\begin{tikzpicture}[node distance=0.32cm, every node/.style={font=\footnotesize, align=center},
  caja/.style={draw=black!60, rounded corners=2pt, minimum height=1.25cm, text width=2.05cm, fill=black!3},
  fuera/.style={caja, draw=red!50!black, fill=red!5},
  queda/.style={caja, draw=blue!60!black, fill=blue!6, text width=2.5cm},
  flecha/.style={-{Stealth[length=2mm]}, black!60}]
\node[caja] (p) {Brecha A--B\\de 1.44\,t/ha\\¿causa del lugar?};
\node[fuera, right=of p] (c) {Clima y\\altitud\\Sec.~\ref{sec:clima}};
\node[fuera, right=of c] (f) {Fenología\\y verdor\\Sec.~\ref{sec:fenologia}};
\node[fuera, right=of f] (s) {Suelo\\(INEGI)\\Sec.~\ref{sec:suelo}};
\node[fuera, right=of s] (v) {Ceniza del\\Popocatépetl\\Sec.~\ref{sec:volcan}};
\node[fuera, right=of v] (o) {Nivel oficial\\del municipio\\Sec.~\ref{sec:siap}};
\node[queda, right=of o] (q) {Forma de obtener\\la etiqueta en\\cada municipio};
\draw[flecha] (p) -- (c); \draw[flecha] (c) -- (f); \draw[flecha] (f) -- (s); \draw[flecha] (s) -- (v);
\draw[flecha] (v) -- (o); \draw[flecha] (o) -- (q);
\end{tikzpicture}
\caption{Recorrido del análisis. En rojo, las causas de lugar que probamos y descartamos. En azul, la explicación que
queda en pie.}
\label{fig:embudo}
\end{figure*}
